% !TEX root = ../main.tex

\section{Introduction}

Hierarchical clustering is attractive when the analyst wants more than a flat
partition. A tree exposes coarse-to-fine structure, keeps near-miss cluster
boundaries visible, and provides a natural substrate for interpretable cluster
decisions. In practice, however, most hierarchical
workflows still cut the tree using geometric heuristics such as a fixed height,
an inconsistency threshold, or a manually chosen number of clusters. Those
choices are often difficult to defend when the goal is not only descriptive
grouping but also reproducible statistical claims.

KL-TE takes a different view of the tree. An internal node is not merely a
merge artifact; it is a decision point. For a parent node $u$ with children
$l(u)$ and $r(u)$, the core question is: should this split survive in the final
partition? A linkage tree alone cannot answer that question. It proposes
candidate splits, but it does not tell us whether at least one child genuinely
departs from the parent's feature distribution, whether the two children are
actually distinct enough to represent separate clusters, or whether scanning
the entire tree has produced false positives.

The method therefore requires two kinds of evidence. First, a split should not
survive unless at least one child meaningfully departs from its ancestral
context. Second, even that is not enough: a parent can contain one drifting
branch or a mild asymmetry without supporting two stable child clusters, so the
children themselves must also be distinguishable from one another. This is why
the method uses both a child-parent edge test and a sibling test. The final
partition keeps a split only when both questions are answered favorably after
multiple-testing correction.

Four inferential complications shape the design. The child is a subset of its
parent, so the edge comparison is nested rather than an ordinary two-sample
test. The edge evidence is high-dimensional and correlated, so naive quadratic
aggregation would overcount redundant or noisy directions. Many candidate edges
are examined across the tree, so valid local $p$-values still require
multiplicity control. Finally, sibling comparisons are evaluated on a hierarchy
constructed from the same data, so the raw sibling statistics are selected and
inflated. The implemented method should be read as one response to these four
problems.

The implementation works on a discrete feature representation. Binary data
enter directly. Categorical data are represented through indicator-expanded
binary features. Continuous data are discretized before the same tree-and-test
pipeline is applied. Once the representation is fixed, the inferential objects
are no longer individual samples but subtrees, so each node carries a
feature-wise distribution estimated from its descendant leaves. Statistical
test annotations are computed before the actual tree walk that creates cluster
assignments. This separation between evidence calculation and final traversal
is useful both computationally and scientifically: the statistical support for
each potential split is stored explicitly and can be audited independently of
the resulting partition.

The manuscript follows that decision problem in order. It first constructs the
tree and node distributions that define the candidate splits. It then asks
whether a child departs from its parent, explains why that test needs nested
variance and spectral whitening, and applies hierarchical multiple-testing
control across the tree. It next asks whether the two children under the same
parent are genuinely distinct, and explains why that sibling comparison needs a
post-selection calibration step. Only after both evidence tables are in place
does the final tree traversal convert them into clusters.

At this stage the writing goal is to capture the statistical structure of the
method cleanly enough that the later experimental section can be attached
without rethinking the exposition. The immediate next step after this draft is
to add calibration and power studies that quantify when the inferential claims
made by these tests are empirically trustworthy.
